\chapter{Data, Questions, and the ML Workflow}\label{chap:data-workflow}
\courseweek{Weeks 1--2.} A useful ML system begins before model selection. It begins with a decision, a measurable target, a representation of the available evidence, and a procedure that can be repeated. This chapter establishes the working vocabulary used throughout the course.

\begin{learningbox}
\begin{itemize}
  \item Distinguish a decision problem from an ML task.
  \item Identify an observation, feature vector $x$, target $y$, and prediction $\hat{y}$.
  \item Create a small, inspectable data workflow in Python.
  \item State what a plot or summary statistic does \emph{not} establish.
\end{itemize}
\end{learningbox}

\section{From an engineering question to a learning task}
An engineering question such as ``Which process condition is likely to produce a defect?'' must be made operational. Define the unit of observation, the time at which the prediction is made, the target, the available features, and the decision that follows. A classification task predicts a discrete label; a regression task predicts a continuous value. Neither label is meaningful without a stated measurement process and error cost.

Let $\mathcal{D}=\{(x_i,y_i)\}_{i=1}^{n}$ denote observed examples. A learning algorithm selects a function $f_\theta$ so that $\hat{y}=f_\theta(x)$ is useful on data not used to choose $\theta$. The course treats this as an evidence problem: the data split and evaluation protocol are part of the model, not afterthoughts \citep{mitchell1997}.

\section{A reproducible first pass}
Before fitting a model, inspect shape, column meaning, units, missingness, and plausible ranges. The following pattern is deliberately small enough to read line by line.

\begin{lstlisting}
import numpy as np
import pandas as pd
import matplotlib.pyplot as plt

rng = np.random.default_rng(2026)
df = pd.DataFrame({
    "temperature_c": rng.normal(72, 5, 40),
    "pressure_kpa": rng.normal(210, 18, 40),
})
df["defect_flag"] = (
    (df["temperature_c"] > 75) & (df["pressure_kpa"] > 220)
).astype(int)

print(df.shape)
print(df.describe())
plt.scatter(df["temperature_c"], df["pressure_kpa"],
            c=df["defect_flag"], cmap="coolwarm")
plt.xlabel("Temperature (C)")
plt.ylabel("Pressure (kPa)")
plt.show()
\end{lstlisting}

A coloured scatter plot can reveal a possible association. It does not prove causality, establish a production threshold, or show that a future batch will behave similarly. The lab asks students to write both the supported observation and a limitation.

\begin{labbox}
\textbf{Week 1:} \colab{week-01-lab.ipynb}. Inspect a measurement table, compute descriptive statistics, create one figure, and write an evidence statement. The target is not yet a high-scoring model; it is a clear workflow that another person can reproduce.
\end{labbox}

\section{Data contracts}
A data contract records the feature names, units, allowable values, target definition, split rule, and known missingness. It is lightweight documentation with a strong engineering benefit: it prevents silent changes from being mistaken for model improvement. For every course notebook, include a first cell that reports package versions, a random seed, data provenance, and the shape of the loaded table.

\begin{checkpointbox}
For a proposed predictor of fatigue life, name: (1) one observation, (2) three candidate features, (3) the target and its unit, (4) a decision that uses the prediction, and (5) one feature that would be unavailable at prediction time. Explain why the last feature must be excluded.
\end{checkpointbox}
